% GENERATED FILE. Edit canonical lessons, then run npm run generate:books.
% canonical-source-hash: fnv1a64:4cb8652febfb8f26
% canonical-lessons: LA-C433-iumentum, LA-C433-caballus, LA-C433-mula, LA-C433-clitellae, LA-C433-onero

\chapter{Beast of burden, Nag, She-mule, Pack saddle, To load}
\label{ch:la-beast-of-burden-nag-she-mule-pack-saddle-to-load}
\hlchaptermodality{433}

\begin{chapteropening}
\textbf{By the end of this chapter:} \emph{I can talk about animals that carry loads: a beast of burden, a nag, a she-mule, a pack saddle, and loading them.}

Complete the last lesson of chapter 433: I can talk about animals that carry loads: a beast of burden, a nag, a she-mule, a pack saddle, and loading them.
\end{chapteropening}

\section[iūmentum, iūmentī]{iūmentum, iūmentī — a beast of burden, a draught animal}

\label{lesson:LA-C433-iumentum}



\begin{quote}
Before the new one: say the Latin for to jump down, then the Latin for to kick.
\end{quote}

\subsection*{You'll want to know: iūmentum, iūmentī}
\textbf{iūmentum, iūmentī} — \textquotedblleft{}a beast of burden, a draught animal\textquotedblright{}.

\begin{quote}
Iūmenta sarcinās mercātōrum per montēs portant.
\end{quote}

\begin{quote}
\emph{Pack animals carry the merchants' bags over the mountains.}
\end{quote}

Second declension, neuter: \emph{iūmentum, iūmentī n.} Any animal that pulls or carries: horse, mule, ass or ox.

\begin{cousinweb}
\emph{iūmentum} is usually traced to \emph{iungō} 'to yoke' (Lewis \& Short, from an older \emph{iugimentum}): an animal put under the yoke. EtymDB lists it without a parent, so the line is hedged.
\end{cousinweb}

\begin{grammarlens}[title={What you've built}]
One more word for telling, explaining and describing at B1.
\end{grammarlens}

\subsection*{Your turn}
\begin{itemize}
\raggedright
  \item \emph{Say it:} \emph{iūmentum, iūmentī}
  \item \emph{Say it:} \emph{iūmentum, iūmentī}, once more
  \item \emph{Say it:} a sentence of your own with \emph{iūmentum, iūmentī}
  \item \emph{Recall:} say the Latin for to jump down, then the Latin for to kick, then say \emph{iūmentum, iūmentī} again
\end{itemize}

\begin{tcolorbox}[breakable,colback=teal!4,colframe=teal!35!black,title={Before you move on}]
What does iūmentum, iūmentī mean? (\textquotedblleft{}A beast of burden, a draught animal\textquotedblright{}.) Say it once more.
\end{tcolorbox}
\section[caballus, caballī]{caballus, caballī — a nag, a pack horse}

\label{lesson:LA-C433-caballus}



\begin{quote}
Before the new one: say the Latin for to kick, then the Latin for a beast of burden.
\end{quote}

\subsection*{You'll want to know: caballus, caballī}
\textbf{caballus, caballī} — \textquotedblleft{}a nag, a pack horse\textquotedblright{}.

\begin{quote}
Caballus vetus et lentus est, sed sarcinās portat.
\end{quote}

\begin{quote}
\emph{The nag is old and slow, but it carries the bags.}
\end{quote}

Second declension, masculine: \emph{caballus, caballī m.} An everyday, low word; the fine word is \emph{equus}. French \emph{cheval} and English \emph{cavalry} come from it.

\begin{cousinweb}
\emph{caballus} is usually traced to Gaulish \emph{caballos} 'pack horse' (EtymDB); Lewis \& Short also call it perhaps Celtic and compare Greek \emph{kaballēs}. French \emph{cheval} and, through Italian, English \emph{cavalry} go back to it.
\end{cousinweb}

\begin{grammarlens}[title={What you've built}]
One more word for telling, explaining and describing at B1.
\end{grammarlens}

\subsection*{Your turn}
\begin{itemize}
\raggedright
  \item \emph{Say it:} \emph{caballus, caballī}
  \item \emph{Say it:} \emph{caballus, caballī}, once more
  \item \emph{Say it:} a sentence of your own with \emph{caballus, caballī}
  \item \emph{Recall:} say the Latin for to kick, then the Latin for a beast of burden, then say \emph{caballus, caballī} again
\end{itemize}

\begin{tcolorbox}[breakable,colback=teal!4,colframe=teal!35!black,title={Before you move on}]
What does caballus, caballī mean? (\textquotedblleft{}A nag, a pack horse\textquotedblright{}.) Say it once more.
\end{tcolorbox}
\section[mūla, mūlae]{mūla, mūlae — a she-mule}

\label{lesson:LA-C433-mula}



\begin{quote}
Before the new one: say the Latin for a beast of burden, then the Latin for a nag.
\end{quote}

\subsection*{You'll want to know: mūla, mūlae}
\textbf{mūla, mūlae} — \textquotedblleft{}a she-mule\textquotedblright{}.

\begin{quote}
Mūla parva est sed fortis.
\end{quote}

\begin{quote}
\emph{The she-mule is small but strong.}
\end{quote}

First declension, feminine: \emph{mūla, mūlae f.}, the female of \emph{mūlus}. Romans said a mule never foals, so \emph{cum mūla pepererit} meant 'never'.

\begin{cousinweb}
\emph{mūla} is the feminine of \emph{mūlus} 'mule' (Lewis \& Short). EtymDB traces \emph{mūlus} to PIE \emph{mukslós}.
\end{cousinweb}

\begin{grammarlens}[title={What you've built}]
One more word for telling, explaining and describing at B1.
\end{grammarlens}

\subsection*{Your turn}
\begin{itemize}
\raggedright
  \item \emph{Say it:} \emph{mūla, mūlae}
  \item \emph{Say it:} \emph{mūla, mūlae}, once more
  \item \emph{Say it:} a sentence of your own with \emph{mūla, mūlae}
  \item \emph{Recall:} say the Latin for a beast of burden, then the Latin for a nag, then say \emph{mūla, mūlae} again
\end{itemize}

\begin{tcolorbox}[breakable,colback=teal!4,colframe=teal!35!black,title={Before you move on}]
What does mūla, mūlae mean? (\textquotedblleft{}A she-mule\textquotedblright{}.) Say it once more.
\end{tcolorbox}
\section[clītellae, clītellārum]{clītellae, clītellārum — a pack saddle (for an ass or mule)}

\label{lesson:LA-C433-clitellae}



\begin{quote}
Before the new one: say the Latin for a nag, then the Latin for a she-mule.
\end{quote}

\subsection*{You'll want to know: clītellae, clītellārum}
\textbf{clītellae, clītellārum} — \textquotedblleft{}a pack saddle (for an ass or mule)\textquotedblright{}.

\begin{quote}
Asinus clītellās portat et in clītellīs saccī sunt.
\end{quote}

\begin{quote}
\emph{The donkey carries a pack saddle, and on the pack saddle there are sacks.}
\end{quote}

First declension, feminine plural: \emph{clītellae, clītellārum f.} One thing, plural form.

\begin{cousinweb}
\emph{clītellae} is inherited: EtymDB traces it to Proto-Italic \emph{kleitrelā}. Lewis \& Short give the sense but no origin.
\end{cousinweb}

\begin{grammarlens}[title={What you've built}]
One more word for telling, explaining and describing at B1.
\end{grammarlens}

\subsection*{Your turn}
\begin{itemize}
\raggedright
  \item \emph{Say it:} \emph{clītellae, clītellārum}
  \item \emph{Say it:} \emph{clītellae, clītellārum}, once more
  \item \emph{Say it:} a sentence of your own with \emph{clītellae, clītellārum}
  \item \emph{Recall:} say the Latin for a nag, then the Latin for a she-mule, then say \emph{clītellae, clītellārum} again
\end{itemize}

\begin{tcolorbox}[breakable,colback=teal!4,colframe=teal!35!black,title={Before you move on}]
What does clītellae, clītellārum mean? (\textquotedblleft{}A pack saddle (for an ass or mule)\textquotedblright{}.) Say it once more.
\end{tcolorbox}
\section[onerō, onerāre]{onerō, onerāre — to load, to burden}

\label{lesson:LA-C433-onero}



\begin{quote}
Before the new one: say the Latin for a she-mule, then the Latin for a pack saddle.
\end{quote}

\subsection*{You'll want to know: onerō, onerāre}
\textbf{onerō, onerāre} — \textquotedblleft{}to load, to burden\textquotedblright{}.

\begin{quote}
Servī mūlōs sarcinīs onerant.
\end{quote}

\begin{quote}
\emph{The slaves load the mules with baggage.}
\end{quote}

\emph{onerō, onerāre, onerāvī, onerātum}, with the thing in the ablative. English \emph{onerous} and \emph{exonerate} come from the same root.

\begin{cousinweb}
\emph{onerō} is built on \emph{onus, oneris} 'a load, a burden', from PIE \emph{h\textsubscript{3}énh\textsubscript{2}os} (EtymDB; Lewis \& Short). English \emph{onerous} and \emph{exonerate} come from the same family.
\end{cousinweb}

\begin{grammarlens}[title={What you've built}]
One more word for telling, explaining and describing at B1.
\end{grammarlens}

\subsection*{Your turn}
\begin{itemize}
\raggedright
  \item \emph{Say it:} \emph{onerō, onerāre}
  \item \emph{Say it:} \emph{onerō, onerāre}, once more
  \item \emph{Say it:} a sentence of your own with \emph{onerō, onerāre}
  \item \emph{Recall:} say the Latin for a she-mule, then the Latin for a pack saddle, then say \emph{onerō, onerāre} again
\end{itemize}

\begin{tcolorbox}[breakable,colback=teal!4,colframe=teal!35!black,title={Before you move on}]
What does onerō, onerāre mean? (\textquotedblleft{}To load, to burden\textquotedblright{}.) Say it once more.
\end{tcolorbox}
